% ID: SEL-D0048-OB-E04
% ESTADO: cotejada. Segunda pasada visual y borrador compilado revisados.
% Fuente: Selectividad/sel_2013_matematicas_aplicadas/junio/reserva a/EXAMEN 5 -ANDALUCÍA.pdf
% SHA-256: 76b6f6a138ec0917a2e3b3a2925bccca739af3363e419603b6814d835860441e
% Página: 2; opción B; ejercicio original 4.

En una población próxima a un puerto deportivo se quiere estimar la proporción de habitantes que navegan al menos una vez a la semana. Se toma una muestra, al azar, de 400 habitantes de la población, de los que 160 afirman navegar al menos una vez en semana.
\begin{enumerate}[label=\alph*)]
\item \textbf{(1.5 puntos)} Halle el intervalo de confianza del 90\% para la proporción de habitantes que navegan al menos una vez en semana.
\item \textbf{(1 punto)} A la vista del resultado, se pretende repetir la experiencia para conseguir una cota del error de 0.1 con el mismo nivel de confianza del apartado anterior. \textquestiondown{}Cuántos individuos debe tener al menos la muestra?
\end{enumerate}
